% ABGELOEST am 2026-09-09: gekuerzt als Subsection "Multireference
% geometries: attempted, out of scope" (sec:mr-whynotmr) in
% docs/section_mr_outofscope.tex (Outlook-Kapitel). Nur noch Verlauf.
% "Feasibility of a multireference reference", aus dem MR-Kapitel der Thesis,
% hier abgelegt am 2026-09-07. Wird laut Autor auf einen Absatz im
% OMol25-Kapitel schrumpfen; der Background verweist darauf.
% Referenzierte Labels, die hier nicht definiert sind: sec:mr-setup,
% sec:mr-together, sec:mr-fix.

\section{Feasibility of a multireference reference}
\label{sec:mr-whynotmr}

A broken-symmetry reference is a pragmatic construction: a single determinant
that is not a spin eigenfunction. The alternative is a multireference treatment,
which describes the separation of the electron pair without spin contamination.
This route was attempted for these reactions and abandoned. The reasons are
reported because they bear on what a benchmark of this kind can use.

\paragraph{The blocker is active-space consistency, not gradients.} A CASSCF
transition-state optimisation needs an active space, and the surface being
differentiated is defined by that choice, so the space must remain the same
along the path. Selecting it with AVAS~\cite{sayfutyarova2017} at the DFT
transition state and re-selecting it at the resulting CASSCF geometry returned a
different space, which means the geometry had been optimised with a space that
does not belong to it. The magnitude was measured directly: at rxn7949, changing
the active space at fixed geometry shifted the barrier by 3310\,meV. Eleven of
26 optimisations converged cleanly.

\paragraph{The diagnostic inherits the same ambiguity.} The $M$
diagnostic~\cite{tishchenko2008} would be the natural quantity for stratifying
by multireference character, but it is computed from a CASSCF wavefunction. Two
converged AVAS solutions at the same geometry gave $M = 0.79$ and $M = 0.07$ for
the same reaction, and $M$ halved across geometry changes below 0.3\,\AA. It was
therefore not used for the stratification of Section~\ref{sec:mr-setup};
$N_\mathrm{FOD}$, computed from a single DFT calculation at a fixed geometry,
was used instead.

\fbox{\begin{minipage}{0.95\linewidth}
\textbf{Key takeaway.} The obstacle to a multireference reference at this system
size is not the availability of gradients. It is that the active space --- and
with it the diagnostic that would indicate whether one is needed --- is not
stable against the choice of solution or against small changes in geometry.
\end{minipage}}

\medskip
\paragraph{Automated selection does not remove the problem at this size.}
autoCAS~\cite{stein2019autocas} prescribes the required iteration: re-select the
active space in the optimised orbital basis until the space selected at the
converged geometry is the one used to optimise it. It requires a DMRG
entanglement calculation over the full valence orbital space, which is not
practical at 11--12 atoms with this basis set. Protocols that enforce
consistency across geometries exist~\cite{seal2025wasp, bensberg2023}, but they
assume a fixed active space; adapting its composition along a reaction
coordinate is stated there as open work.

\paragraph{Requirements for a future attempt.} Three routes remain. A
self-consistency iteration at the CASSCF level, repeated until the selected
space is a fixed point of the optimisation, then enlarged to check that no
relevant orbital was omitted. A move to CASPT2 or NEVPT2 geometries, which add
the dynamic correlation CASSCF omits --- CASSCF accentuates single and double
bond character, which at a transition state means overstating diradical
character, and bond lengths differ from CASPT2 by 0.01 to 0.02\,\AA{} for small
organic molecules~\cite{page2003}; analytic gradients for these
methods~\cite{park2019} are not available in the software stack used here, and
numerical gradients cost of order $6N$ single points per optimisation step. Or a
reference set restricted to systems small enough for near-exact methods, which
will not resemble the reactions of interest.

The consequence for this chapter is that the broken-symmetry surface is used
with its limitations stated. One observation supports the choice after the fact:
where several methods find a saddle point on that surface, they find the same
one (Sections~\ref{sec:mr-together} and~\ref{sec:mr-fix}); the surface has a
reproducible stationary point that independent searches converge to.
